%%%%%%%%%%%%%%%%%%%%%%%%%%%%%%%%%%%%%%%%%%%%
%             DOCUMENT SETTINGS
%%%%%%%%%%%%%%%%%%%%%%%%%%%%%%%%%%%%%%%%%%%%

\setlength{\parindent}{2em}
\setlength{\parskip}{5pt}
\newlength{\originalparindent}
\newlength{\originalparskip}
\AtBeginDocument{%
    \setlength{\originalparindent}{\parindent}%
    \setlength{\originalparskip}{\parskip}%
}

% A5 PAPER
\usepackage[a5paper,top=0.75in,bottom=0.75in,left=0.75in,right=0.75in,marginparwidth=1.75cm]{geometry}

% SECTION TITLE FONT SIZES (move here from preset)
\titleformat*{\section}{\LARGE\bfseries}
\titleformat*{\subsection}{\Large\bfseries}
\titleformat*{\subsubsection}{\large\bfseries}
\titleformat*{\subsubsubsection}{\large\bfseries}

% CUSTOM SECTION COMMANDS
\newcommand{\module}[1]{%
    \clearpage
    \refstepcounter{section}
    \vspace*{0.5cm}
    {\LARGE\bfseries\noindent Module \thesection\par}
    \vspace{0.5cm}
    {\noindent\Huge\bfseries #1\par}
    \vspace{1.5cm}
    \addcontentsline{toc}{section}{Module \thesection : #1}
}

\newcommand{\unitsec}[1]{%
    \clearpage
    \refstepcounter{section}
    \vspace*{0.5cm}
    {\LARGE\bfseries\noindent Unit \thesection\par}
    \vspace{0.5cm}
    {\noindent\huge\bfseries #1\par}
    \vspace{1.5cm}
    \addcontentsline{toc}{section}{Unit \thesection : #1}
}

\newcommand{\chaptersec}[1]{%
    \clearpage
    \refstepcounter{section}
    \vspace*{0.5cm}
    {\LARGE\bfseries\noindent Chapter \thesection\par}
    \vspace{0.5cm}
    {\noindent\huge\bfseries #1\par}
    \vspace{1.5cm}
    \addcontentsline{toc}{section}{Chapter \thesection : #1}
}

%%% PROBLEM ENVIRONMENT (was in preset.tex, now here)
\newenvironment{problem}{%
    \hrule
    \vspace{5pt}
}{%
    \vspace{10pt}\hrule\vspace{10pt}
}

\newenvironment{solution_pset}{%
    \noindent\textbf{\Large Solution}
}{%
    \newpage
}

\numberwithin{equation}{subsection}

%%% DEFINITION ENVIRONMENT
% Definition counter linked to section
\newcounter{definition}[section]
\renewcommand{\thedefinition}{\thesection.\arabic{definition}}

\newtcolorbox{definition}[2][]{%
    colback=white,
    colframe=black,
    boxrule=0.5pt,
    arc=0pt,
    left=6pt,
    right=6pt,
    top=6pt,
    bottom=6pt,
    fonttitle=\bfseries,
    title={%
        \stepcounter{definition}% Increment counter
        Definition \thedefinition{} : #2% Then display
    },
    #1
}

%%% EXAMPLE ENVIRONMENT
% Example counter linked to section
\newcounter{example}[section]
\renewcommand{\theexample}{\thesection.\arabic{example}}

\newtcolorbox{example}[2][]{%
    breakable,
    colback=white,
    colframe=blue,
    boxrule=0.5pt,
    arc=4pt,
    left=6pt,
    right=6pt,
    top=6pt,
    bottom=6pt,
    fonttitle=\bfseries,
    parbox=false,  % Don't use parbox mode
    before upper={%
        \setlength{\parindent}{\originalparindent}%
        \setlength{\parskip}{\originalparskip}%
        \noindent%
    },
    title={%
        \refstepcounter{example}%
        Example \theexample{} : #2%
    },
    #1
}

%%% THEOREM ENVIRONMENT
% Theorem counter linked to section
\newcounter{theorem}[section]
\renewcommand{\thetheorem}{\thesection.\arabic{theorem}}
\definecolor{darkred}{RGB}{139, 0, 0} 
\newenvironment{theorem}[2][]{%
    \refstepcounter{theorem}%
    \begin{tcolorbox}[
        colback=white,
        colframe=darkred,
        coltitle=white,
        colbacktitle=darkred,
        boxrule=0.5pt,
        arc=0pt,
        left=6pt,
        right=6pt,
        top=6pt,
        bottom=6pt,
        fonttitle=\bfseries,
        title={Theorem \thetheorem{} : #2},
        #1
    ]
}{%
    \end{tcolorbox}%
}

%%% NOTE ENVIRONMENT
\definecolor{darkgray}{RGB}{100, 100, 100}

\newtcolorbox{note}{%
    colback=white,
    colframe=darkgray,
    boxrule=0.5pt,
    arc=0pt,
    left=6pt,
    right=6pt,
    top=6pt,
    bottom=6pt,
    before upper={%
        \noindent{\textit{Note : }}\par\vspace{1em}%
    }
}

%%% REMARK ENVIRONMENT
\newtcolorbox{remark}{%
    colback=white,
    colframe=darkgray,
    boxrule=0.5pt,
    arc=0pt,
    left=6pt,
    right=6pt,
    top=6pt,
    bottom=6pt,
    before upper={%
        \noindent\textbf{\textit{Remark : }}\par\vspace{1em}%
    }
}
